\documentclass[11pt]{article}

%-----------------------------------------------------------------------------%
% PACKAGES
%-----------------------------------------------------------------------------%
\usepackage[margin=1in]{geometry}
\usepackage{amsmath,amssymb,amsthm}
\usepackage{mathtools}
\usepackage{bm}
\usepackage{hyperref}
\usepackage{enumitem}
\usepackage{graphicx}
\usepackage{booktabs}
\usepackage{mathrsfs}
\usepackage{microtype}
\usepackage{authblk}

\hypersetup{
  colorlinks=true,
  linkcolor=blue,
  citecolor=blue,
  urlcolor=blue
}

%-----------------------------------------------------------------------------%
% THEOREM ENVIRONMENTS
%-----------------------------------------------------------------------------%
\newtheorem{definition}{Definition}[section]
\newtheorem{theorem}[definition]{Theorem}
\newtheorem{proposition}[definition]{Proposition}
\newtheorem{corollary}[definition]{Corollary}
\newtheorem{remark}[definition]{Remark}
\newtheorem{lemma}[definition]{Lemma}

%-----------------------------------------------------------------------------%
% MACROS
%-----------------------------------------------------------------------------%
\newcommand{\Z}{\mathbb{Z}}
\newcommand{\N}{\mathbb{N}}
\newcommand{\R}{\mathbb{R}}
\newcommand{\F}{\mathbb{F}}
\newcommand{\G}{\mathbb{G}}
\newcommand{\calS}{\mathcal{S}}
\newcommand{\calT}{\mathcal{T}}
\newcommand{\calM}{\mathcal{M}}
\newcommand{\calK}{\mathcal{K}}
\newcommand{\calE}{\mathcal{E}}
\newcommand{\calF}{\mathcal{F}}

\newcommand{\Enc}{\mathsf{Enc}}
\newcommand{\Dec}{\mathsf{Dec}}

\DeclareMathOperator{\EncOp}{E}
\DeclareMathOperator{\HKDF}{HKDF}
\DeclareMathOperator{\HMAC}{HMAC}
\DeclareMathOperator{\SHA}{SHA256}

\newcommand{\norm}[1]{\left\lVert #1 \right\rVert}

%-----------------------------------------------------------------------------%
% TITLE
%-----------------------------------------------------------------------------%

\title{Track B: Multiplicity-Crypto Integrated Certification and Transport\\
Defensive Prior Art Publication and Technical Report}

\author{Citizen Gardens} 
\affil{The Foundation of Multiplicity}

\date{April 27, 2026}

\begin{document}
\maketitle

\begin{abstract}
This document specifies \emph{Track B}, a multiplicity-theoretic certification and transport framework built on top of the native Multiplicity Crypto stack.
Track B binds stateful data products to BN254 Pedersen commitments, prime-indexed multiplicity operators, and transcript-bound HKDF-derived AEAD keys, with optional Quantum-State-Based Security (QSBS) extensions and QKD hybrid key substrates.
The goal of this report is twofold:
(1) to provide a mathematically complete specification suitable for independent re-implementation, and
(2) to constitute prior art for defensive patent purposes, preventing third-party claims on the techniques described herein.
\end{abstract}
\newpage
\tableofcontents
\newpage
%-----------------------------------------------------------------------------%
\section{Executive Summary}
%-----------------------------------------------------------------------------%

\subsection{Purpose of Track B}

Track B is defined as a second, cryptography-focused track in the PhaseMirror / Multiplicity ecosystem whose primary role is to \emph{certify} and \emph{securely transport} evolving data products and model states.
Rather than treating Track B as a simple logging or checksum layer, the design interprets Track B as a multiplicity-governed protocol:
each transformation, message, and acknowledgement becomes a prime-labeled interaction step whose history is encoded and preserved in a multiplicity operator.

Concretely, Track B integrates:
\begin{itemize}[leftmargin=1.5em]
  \item A BN254-based Pedersen commitment engine (via Multiplicity Crypto) for circuit-native commitments.
  \item A prime-indexed multiplicity operator that G\"odel-encodes Track B transcripts into uniquely factorizable integers.
  \item A transcript hash chain and context hash used as associated data for authenticated encryption.
  \item An HKDF-based key schedule with multiplicity labels embedded into the info string, producing directional AEAD keys.
  \item An optional QSBS overlay that maps classical payloads and quantum-style amplitudes into a joint frequency space and couples them through a low-rank tensor.
\end{itemize}

\subsection{Defensive Publication Scope}

This document is intended as a defensive prior-art publication and covers at least the following contributions related to Track B:

\begin{enumerate}[leftmargin=1.5em,label=\textbf{B.\arabic*}]
  \item \textbf{Track B state certification} via BN254 Pedersen commitments over frequency-mapped state representations and multiplicity-labeled transcripts.
  \item \textbf{Multiplicity-indexed HKDF derivation for Track B} directional keys, including an 8-byte multiplicity label embedded in the info string.
  \item \textbf{Frequency-unified state commitments for Track B}, where classical and (optionally) quantum-like state descriptors are combined into a single commitment handle.
  \item \textbf{Prime-indexed multiplicity transcript binding for Track B}, giving recursively stable transcript identities across Track B turns.
  \item \textbf{Adaptive Track B key rotation governed by feedback dynamics} on error, latency, and load metrics, expressed in a multiplicity-theoretic state space.
\end{enumerate}

All such techniques are hereby described with sufficient detail for reconstruction and are intended to be irrevocably placed into the public domain for defensive purposes.

%-----------------------------------------------------------------------------%
\section{Track B Conceptual Architecture}
%-----------------------------------------------------------------------------%

\subsection{Layered View}

Track B is organized as a layered system, specialized for certification and transport:

\begin{itemize}[leftmargin=1.5em]
  \item \textbf{B0 -- Curve and primitives:} BN254 group and scalar field, random oracle, AEAD scheme.
  \item \textbf{B1 -- Commitment layer:} BN254 Pedersen commitments for Track B states.
  \item \textbf{B2 -- Frequency mapping:} maps Track B states into joint frequency vectors.
  \item \textbf{B3 -- Multiplicity and transcripts:} prime-indexed transcript encoding and multiplicity operator.
  \item \textbf{B4 -- Key derivation:} HKDF-based directional keys with multiplicity labels.
  \item \textbf{B5 -- Encryption/AEAD:} AEAD encryption with transcript-bound associated data.
  \item \textbf{B6 -- Feedback dynamics:} adaptive key rotation and state evolution controlled by observed metrics.
  \item \textbf{B7 -- Optional QSBS overlay:} joint classical/quantum frequency model integrated into the above.
\end{itemize}

Each layer is designed to be independently testable but composable, with clear functional boundaries and simple interfaces.

\subsection{Track B State and Roles}

Track B is parameterized by a discrete time index \(t \in \N\) and involves at least two roles (e.g., producer and consumer, or module A and module B).
Its instantaneous state is:

\[
  S_t^{(B)} = \bigl(x_t^{(B)}, \sigma_t^{(B)}, M_t^{(B)}, T_t^{(B)}\bigr),
\]

optionally extended with a quantum or amplitude-like component \(|\psi_t^{(B)}\rangle\) when QSBS is used.

\begin{itemize}[leftmargin=1.5em]
  \item \(x_t^{(B)}\): classical Track B payload (bitstring or structured data serialized to bits).
  \item \(\sigma_t^{(B)}\): Track B transcript (ordered message history and labels).
  \item \(M_t^{(B)}\): multiplicity operator, a vector of prime-indexed message-class counts.
  \item \(T_t^{(B)}\): coupling tensor, modeling interactions between classical and quantum modes.
\end{itemize}

The goal is that every emission and transformation of Track B is both cryptographically committed and multiplicity-identifiable.

%-----------------------------------------------------------------------------%
\section{Mathematical Overview}
%-----------------------------------------------------------------------------%

\subsection{Standing Assumptions}

\paragraph{Cryptographic primitives.}
We assume:
\begin{itemize}[leftmargin=1.5em]
  \item A pairing-friendly elliptic curve BN254 with group \(\G_1\) of prime order \(r\) and generator \(G\).
  \item A random-oracle hash function \(\SHA: \{0,1\}^* \to \{0,1\}^{256}\).
  \item An HMAC-based HKDF key derivation function.
  \item A nonce-based AEAD scheme (e.g., AES-256-GCM) that is IND-CPA and INT-CTXT secure.
\end{itemize}

\paragraph{Security parameter.}
Let \(\lambda\) denote the security parameter; all polynomial-time bounds are with respect to \(\lambda\).
Negligible functions are denoted \(\mathsf{negl}(\lambda)\).

\subsection{Track B Pedersen Commitments}

\begin{definition}[BN254 Pedersen commitment for Track B]
Let \(\G_1\) be the prime-order subgroup of BN254 with generator \(G\), and \(\F_r\) the scalar field of order \(r\).
Define an independent generator \(H \in \G_1\) as
\[
  H = h_G \cdot G, \quad
  h_G = \bigl(\mathrm{int}(\SHA(\texttt{"pedersen\_h"} \| \mathrm{enc}(G))) \bmod r\bigr) \in \F_r.
\]
For a Track B value \(v \in \F_r\) and blinding factor \(\rho \in \F_r\), the commitment is
\[
  C^{(B)}(v, \rho) = \rho G + v H \in \G_1.
\]
\end{definition}

This is a standard Pedersen commitment specialized to BN254 and used as the canonical commitment of Track B state encodings.

\subsection{Prime-Indexed Transcript Encoding and Multiplicity}

Let \(\{p_i\}_{i \ge 1}\) denote the increasing sequence of primes.

\begin{definition}[Track B message classes]
Let \(\mathcal{C} = \{c_1, \dots, c_k\}\) be the set of distinct Track B message classes (e.g., \texttt{state\_update}, \texttt{ack}, \texttt{error}, etc.).
For a transcript \(\sigma_t^{(B)}\), let \(c_i(\sigma_t^{(B)}) \in \N\) be the count of messages of class \(c_i\) up to time \(t\).
\end{definition}

\begin{definition}[Prime-indexed transcript encoding]
The Track B transcript encoding is the integer
\[
  \mathrm{enc}_B(\sigma_t^{(B)}) = \prod_{i=1}^k p_i^{c_i(\sigma_t^{(B)})}.
\]
\end{definition}

\begin{definition}[Track B multiplicity operator]
The Track B multiplicity operator at time \(t\) is the vector
\[
  M_t^{(B)} = \bigl(c_1(\sigma_t^{(B)}), \dots, c_k(\sigma_t^{(B)})\bigr) \in \Z_{\ge 0}^k.
\]
\end{definition}

\begin{proposition}[Injectivity of Track B transcript encoding]
Distinct Track B message-class multisets yield distinct encodings.
Equivalently, the map
\(
\sigma_t^{(B)} \mapsto \mathrm{enc}_B(\sigma_t^{(B)})
\)
is injective on multisets of classes.
\end{proposition}

\begin{proof}
By the Fundamental Theorem of Arithmetic, the factorization of a positive integer into primes is unique.
Since each message class \(c_i\) is assigned its own prime \(p_i\), the exponents of \(p_i\) in \(\mathrm{enc}_B(\sigma_t^{(B)})\) are exactly \(c_i(\sigma_t^{(B)})\).
Thus the multiset \(\bigl\{(c_i, c_i(\sigma_t^{(B)}))\bigr\}\) is uniquely determined by the encoding.
\end{proof}

\begin{theorem}[Operator norms of \(M_t^{(B)}\)]
Let \(N_t = \sum_{i=1}^k c_i(\sigma_t^{(B)})\).
Then:
\begin{align*}
  \norm{M_t^{(B)}}_1 &= N_t, \\
  \norm{M_t^{(B)}}_2 &\le N_t, \\
  \norm{M_t^{(B)}}_\infty &\le N_t.
\end{align*}
\end{theorem}

\begin{proof}
The \(\ell_1\) norm of \(M_t^{(B)}\) is \(\sum_i c_i = N_t\).
Since each component is non-negative and bounded by \(N_t\), we have
\(
\norm{M_t^{(B)}}_2^2 = \sum_i c_i^2 \le \left(\max_i c_i\right)\sum_i c_i \le N_t^2
\),
so \(\norm{M_t^{(B)}}_2 \le N_t\).
Similarly, \(\norm{M_t^{(B)}}_\infty = \max_i c_i \le N_t\).
\end{proof}

\subsection{Frequency Mapping for Track B}

Track B uses a frequency map to embed its state into a Euclidean space suitable for unified commitments.

\begin{definition}[Classical Track B frequency map]
Fix two distinct frequencies \(f_0, f_1 \in \R\).
Let \(x_t^{(B)} \in \{0,1\}^n\) be the bitstring encoding of the Track B payload at time \(t\).
Define the classical frequency map
\[
  F_c^{(B)}(x_t^{(B)}) = (f_{c,1}, \dots, f_{c,n}) \in \R^n,
\]
where
\[
  f_{c,i} =
  \begin{cases}
    f_0 & \text{if } (x_t^{(B)})_i = 0, \\
    f_1 & \text{if } (x_t^{(B)})_i = 1.
  \end{cases}
\]
\end{definition}

Optionally, if Track B incorporates quantum-like or probabilistic amplitude vectors, we can define:

\begin{definition}[Quantum-like Track B frequency map]
Let \(|\psi_t^{(B)}\rangle = \sum_{j=1}^m \alpha_j |j\rangle\) with complex amplitudes \(\alpha_j\).
Define
\[
  F_q^{(B)}(|\psi_t^{(B)}\rangle) = 
  \bigl(|\alpha_1|, \arg \alpha_1, \dots, |\alpha_m|, \arg \alpha_m\bigr) \in \R^{2m},
\]
where \(\arg\) is taken in a fixed principal range (e.g., \([0,2\pi)\)).
\end{definition}

\begin{definition}[Joint Track B frequency vector]
The joint frequency representation of Track B state at time \(t\) is
\[
  F_t^{(B)} =
  \begin{cases}
    F_c^{(B)}(x_t^{(B)}) \in \R^n & \text{(purely classical case)},\\[0.5em]
    \bigl(F_c^{(B)}(x_t^{(B)}), F_q^{(B)}(|\psi_t^{(B)}\rangle)\bigr) \in \R^{n+2m} & \text{(QSBS case)}.
  \end{cases}
\]
\end{definition}

\subsection{Coupling Tensor and Encryption Operator}

In the QSBS-enabled version of Track B, a coupling tensor captures interactions between classical and amplitude-mode components.

\begin{definition}[Track B coupling tensor]
Let \(u_t = F_c^{(B)}(x_t^{(B)}) \in \R^n\) and \(v_t = F_q^{(B)}(|\psi_t^{(B)}\rangle) \in \R^{2m}\).
Define the rank-1 coupling tensor
\[
  T_t^{(B)} = u_t v_t^\top \in \R^{n \times 2m}.
\]
\end{definition}

\begin{proposition}[Norms of \(T_t^{(B)}\)]
Let \(B_c = \sup_t \norm{u_t}_2\) and \(B_q = \sup_t \norm{v_t}_2\).
Then:
\begin{align*}
  \norm{T_t^{(B)}}_F &= \norm{u_t}_2 \norm{v_t}_2 \le B_c B_q, \\
  \norm{T_t^{(B)}}_{2\to 2} &= \norm{u_t}_2 \norm{v_t}_2 \le B_c B_q.
\end{align*}
\end{proposition}

\begin{proof}
For a rank-1 matrix \(T = u v^\top\), the Frobenius norm is
\(
\norm{T}_F^2 = \sum_{ij} u_i^2 v_j^2 = \norm{u}_2^2 \norm{v}_2^2
\),
and the only non-zero singular value is \(\sigma_1 = \norm{u}_2\norm{v}_2\).
Thus \(\norm{T}_F = \norm{T}_{2\to 2} = \norm{u}_2\norm{v}_2\).
The bound follows by definition of \(B_c,B_q\).
\end{proof}

\begin{definition}[Track B encryption operator]
Let \(S_t^{(B)} = (x_t^{(B)}, \sigma_t^{(B)}, M_t^{(B)}, T_t^{(B)})\) and let \(F_B\) be a (possibly non-linear) feedback function.
The Track B encryption operator is
\[
  \EncOp_B(t) = M_t^{(B)} \bigl( T_t^{(B)} S_t^{(B)} \bigr) + F_B(S_t^{(B)}).
\]
\end{definition}

The exact action of \(M_t^{(B)}\) on \(T_t^{(B)} S_t^{(B)}\) depends on how the state is embedded into a vector space;
for the purpose of analysis, it suffices to note that the resulting vector norm is bounded in terms of \(N_t, B_c, B_q\) and bounds on \(F_B\).

\begin{proposition}[Norm bound on Track B encryption output]
Assume there exists \(D_S, D_F > 0\) such that \(\norm{S_t^{(B)}}_2 \le D_S\) and \(\norm{F_B(S_t^{(B)})}_2 \le D_F\) for all \(t\).
Then
\[
  \norm{\EncOp_B(t)}_2 \le N_t B_c B_q D_S + D_F.
\]
\end{proposition}

\begin{proof}
Using the standard inequalities:
\[
  \norm{\EncOp_B(t)}_2 \le \norm{M_t^{(B)} T_t^{(B)} S_t^{(B)}}_2 + \norm{F_B(S_t^{(B)})}_2.
\]
Treating \(M_t^{(B)}\) as a diagonal operator with spectral norm bounded by \(N_t\), and using the operator norms for \(T_t^{(B)}\),
\[
  \norm{M_t^{(B)} T_t^{(B)} S_t^{(B)}}_2 \le \norm{M_t^{(B)}}_{2\to 2} \norm{T_t^{(B)}}_{2\to 2} \norm{S_t^{(B)}}_2 \le N_t B_c B_q D_S.
\]
Adding the bound on \(F_B\) yields the stated inequality.
\end{proof}

\subsection{Feedback Dynamics and Contraction}

\begin{definition}[Track B feedback map]
Let \(\Theta_t = (\varepsilon_t, \lambda_t, \mathrm{QBER}_t, \ell_t)\) denote the vector of observed metrics at time \(t\) (e.g., error rate, latency, quantum bit error rate, system load).
The Track B feedback map is
\[
  f_B : (M_t^{(B)}, T_t^{(B)}, \Theta_t) \mapsto (M_{t+1}^{(B)}, T_{t+1}^{(B)}).
\]
\end{definition}

\begin{definition}[Product metric space]
Equip \(\R^k \times \R^{n\times 2m}\) with the norm
\[
  \norm{(M,T)} := \norm{M}_2 + \norm{T}_F.
\]
\end{definition}

\begin{theorem}[Banach contraction for Track B feedback]
Suppose there exist non-negative constants \(L_M, L_T\) such that for all \(M,T,M',T'\),
\begin{align*}
  \norm{f_B(M,T,\Theta)_M - f_B(M',T',\Theta)_M}_2 &\le L_M \norm{(M,T)-(M',T')}, \\
  \norm{f_B(M,T,\Theta)_T - f_B(M',T',\Theta)_T}_F &\le L_T \norm{(M,T)-(M',T')},
\end{align*}
and let \(L = L_M + L_T\).
If \(L < 1\) for the regime of interest, then:
\begin{enumerate}[leftmargin=1.5em]
  \item There exists a unique fixed point \((M_\ast^{(B)}, T_\ast^{(B)})\).
  \item For any initialization \((M_0^{(B)}, T_0^{(B)})\), the iterates converge geometrically:
  \[
    \norm{(M_t^{(B)}, T_t^{(B)}) - (M_\ast^{(B)}, T_\ast^{(B)})} \le L^t \norm{(M_0^{(B)}, T_0^{(B)}) - (M_\ast^{(B)}, T_\ast^{(B)})}.
  \]
\end{enumerate}
\end{theorem}

\begin{proof}
The space \(\R^k \times \R^{n\times 2m}\) equipped with the given norm is complete.
The Lipschitz assumptions imply
\[
  \norm{f_B(M,T,\Theta) - f_B(M',T',\Theta)} \le (L_M + L_T) \norm{(M,T)-(M',T')} = L \norm{(M,T)-(M',T')}.
\]
Thus \(f_B\) is a contraction with constant \(L < 1\), and the Banach Fixed-Point Theorem applies.
\end{proof}

This provides a mathematical basis for eventual stabilisation of multiplicity labels used in Track B HKDF derivation and for designing key-rotation policies.

%-----------------------------------------------------------------------------%
\section{Track B Key Derivation and Encryption}
%-----------------------------------------------------------------------------%

\subsection{Transcript Chains and Context Hash}

Track B uses per-role transcript chains and a combined context hash for binding keys and associated data to the protocol history.

\begin{definition}[Per-role transcript chain]
For each role \(s \in \{A,B\}\), define the chain recursively:
\[
  \mathrm{chain}_{s,0} = 0^{256}, \qquad
  \mathrm{chain}_{s,t} = \SHA\bigl(\mathrm{chain}_{s,t-1} \, \| \, m_{s,t}\bigr),
\]
where \(m_{s,t}\) is the \(t\)-th message from sender \(s\) in Track B.
\end{definition}

\begin{definition}[Track B context hash]
The Track B context hash at time \(t\) is
\[
  \mathrm{ctx}_t^{(B)} = \SHA\bigl(\mathrm{chain}_{A,t} \, \| \, \mathrm{chain}_{B,t}\bigr).
\]
\end{definition}

\subsection{Multiplicity-Indexed HKDF Derivation}

\begin{definition}[Multiplicity label for Track B]
Let \(M_t^{(B)}\) be the multiplicity vector for Track B at time \(t\).
Define an 8-byte encoding
\[
  \mathrm{mult}_t^{(B)} = \mathrm{Enc}_8\bigl(M_t^{(B)}\bigr),
\]
where \(\mathrm{Enc}_8\) is any injective encoding from the relevant range into 8 bytes (e.g., by hashing or truncating a canonical encoding).
\end{definition}

\begin{definition}[Directional Track B AEAD key]
Let \(K_t\) be a shared secret (e.g., from QKD or an existing key-agreement scheme),
let \(\mathrm{salt}_t\) be a salt, and let \(d \in \{A\to B, B\to A\}\) denote direction.
Define the info string
\[
  \mathrm{info}_d^{(B)} = \texttt{"QSBS-v1:enc:"} \, \| \, d \, \| \, \mathrm{ctx}_t^{(B)} \, \| \, \mathrm{mult}_t^{(B)},
\]
and the directional AEAD key
\[
  K_{\mathrm{enc},d,t}^{(B)} = \HKDF\bigl(K_t; \mathrm{salt}_t, \mathrm{info}_d^{(B)}\bigr).
\]
\end{definition}

This design treats the multiplicity label as authenticated meta-information influencing key derivation, rather than as a source of entropy, preserving compatibility with standard HKDF security analyses.

\subsection{Track B AEAD Encryption}

\begin{definition}[Track B ciphertext]
Let \(x_t^{(B)}\) be the plaintext payload, let \(\mathrm{nonce}_{d,t}\) be a nonce,
and let \(\mathrm{ad}_t^{(B)}\) be associated data defined as
\[
  \mathrm{ad}_t^{(B)} = \bigl(\mathrm{ctx}_t^{(B)}, F_t^{(B)}, M_t^{(B)}\bigr)
\]
encoded in some deterministic, unambiguous format.
The Track B ciphertext is
\[
  C_t^{(B)} = \Enc_{K_{\mathrm{enc},d,t}^{(B)}}\bigl(\mathrm{nonce}_{d,t}, x_t^{(B)}, \mathrm{ad}_t^{(B)}\bigr).
\]
\end{definition}

Decryption proceeds by recomputing the same key and associated data from the transcript and multiplicity state, then verifying the AEAD tag.

%-----------------------------------------------------------------------------%
\section{Track B State Certification}
%-----------------------------------------------------------------------------%

\subsection{Canonical Commitment Handle}

\begin{definition}[Track B commitment handle]
Given a Track B state \(S_t^{(B)}\), compute the frequency representation \(F_t^{(B)}\)
and serialize it to a string \(f_t\) via a deterministic encoding (e.g., fixed-precision decimal).
Let \(v_t \in \F_r\) be the reduction of \(\SHA(f_t)\) modulo \(r\) and let \(\rho_t\) be a scalar derived via a domain-separated hash.
Define
\[
  h_t^{(B)} = C^{(B)}(v_t, \rho_t).
\]
\end{definition}

This handle is used as the canonical Track B commitment to the state at time \(t\).
It can be published, stored, or verified in a ZK circuit without revealing the underlying state.

\subsection{Audit and Verification}

Given a claimed Track B state \(S_t^{(B)}\) and blinding \(\rho_t\), verification consists of:
\begin{enumerate}[leftmargin=1.5em]
  \item Recomputing \(F_t^{(B)}\) and its serialized representation \(f_t\).
  \item Computing \(v_t = \mathrm{int}(\SHA(f_t)) \bmod r\).
  \item Computing \(C^{(B)}(v_t,\rho_t)\) and checking equality to the published handle \(h_t^{(B)}\).
\end{enumerate}

This can be implemented either off-chain (e.g., in Rust/TypeScript) or inside a ZK circuit over BN254.

%-----------------------------------------------------------------------------%
\section{Example Code Snippet}
%-----------------------------------------------------------------------------%

This section provides representative (non-executable in this document) code snippets illustrating a possible implementation of the core Track B primitives using a BN254/WASM stack and a TypeScript wrapper.

\subsection{BN254 Pedersen Commitment (Rust)}

\begin{verbatim}
use ark_bn254::{G1Projective, Fr};
use ark_ec::Group;
use ark_ff::{PrimeField, BigInteger};
use sha2::{Sha256, Digest};
use wasm_bindgen::prelude::*;

/// derive_scalar(d, s) = SHA-256(d || s) mod r
fn derive_scalar(domain: &str, data: &str) -> Fr {
    let mut h = Sha256::new();
    h.update(domain.as_bytes());
    h.update(data.as_bytes());
    let bytes = h.finalize();
    Fr::from_be_bytes_mod_order(&bytes)
}

/// H = SHA-256("pedersen_h" || enc(G)) * G
fn pedersen_generator_h() -> G1Projective {
    let g = G1Projective::generator();
    let h_scalar = derive_scalar("pedersen_h", "generator");
    g * h_scalar
}

/// Compute Track B commitment handle over BN254.
/// Returns x-coordinate in hex for circuit interop.
#[wasm_bindgen]
pub fn trackb_commitment(serialized_state: &str, salt: &str) -> String {
    let g = G1Projective::generator();
    let h = pedersen_generator_h();
    // Reduce SHA-256(serialized_state) mod r for v
    let v = derive_scalar("trackb_value", serialized_state);
    // Domain-separated blinding
    let rho = derive_scalar("trackb_blinding", salt);
    let commitment = (g * rho + h * v).into_affine();
    let x_bytes = commitment.x.into_bigint().to_bytes_be();
    format!("0x{}", hex::encode(&x_bytes))
}
\end{verbatim}

\subsection{Transcript and Key Derivation (TypeScript)}

\begin{verbatim}
import { createHash, createHmac } from "crypto";
import init, { trackb_commitment } from "./wasm/multiplicity_crypto_wasm";

type Direction = "A2B" | "B2A";

export class TrackBTranscript {
  private chainA: Buffer;
  private chainB: Buffer;

  constructor() {
    this.chainA = Buffer.alloc(32, 0);
    this.chainB = Buffer.alloc(32, 0);
  }

  update(role: "A" | "B", message: Buffer) {
    const prev = role === "A" ? this.chainA : this.chainB;
    const next = createHash("sha256").update(prev).update(message).digest();
    if (role === "A") this.chainA = next;
    else this.chainB = next;
  }

  contextHashHex(): string {
    const ctx = createHash("sha256")
      .update(this.chainA)
      .update(this.chainB)
      .digest("hex");
    return "0x" + ctx;
  }

  chains(): { chainA: Buffer; chainB: Buffer } {
    return { chainA: Buffer.from(this.chainA), chainB: Buffer.from(this.chainB) };
  }
}

function hkdfExtract(salt: Buffer, ikm: Buffer): Buffer {
  return createHmac("sha256", salt).update(ikm).digest();
}

function hkdfExpand(prk: Buffer, info: Buffer, length: number): Buffer {
  const okm = Buffer.alloc(length);
  let t = Buffer.alloc(0);
  let offset = 0;
  for (let i = 1; offset < length; i++) {
    t = createHmac("sha256", prk)
      .update(t)
      .update(info)
      .update(Buffer.from([i]))
      .digest();
    t.copy(okm, offset);
    offset += t.length;
  }
  return okm.slice(0, length);
}

/// Derive Track B directional key with multiplicity label.
export function deriveTrackBKey(
  sharedKey: Buffer,
  salt: Buffer,
  direction: Direction,
  ctxHex: string,
  multiplicityLabel: Buffer // 8 bytes
): Buffer {
  const info = Buffer.concat([
    Buffer.from("QSBS-v1:enc:"),
    Buffer.from(direction),
    Buffer.from(ctxHex.replace(/^0x/, ""), "hex"),
    multiplicityLabel,
  ]);
  const prk = hkdfExtract(salt, sharedKey);
  return hkdfExpand(prk, info, 32);
}

/// Compute Track B commitment handle (frequency + multiplicity-based salt).
export async function computeTrackBHandle(
  serializedFrequency: string,
  multiplicityIndex: number
): Promise<string> {
  await init(); // ensure WASM is ready
  const salt = `trackb-${multiplicityIndex.toString(16)}`;
  return trackb_commitment(serializedFrequency, salt);
}
\end{verbatim}

These snippets are illustrative and omit error handling, AEAD integration, and production hardening, but they show how Track B can be concretely wired into a BN254/WASM + TypeScript environment.

%-----------------------------------------------------------------------------%
\section{Implementation Guidelines and Recommendations}
%-----------------------------------------------------------------------------%

\subsection{Track B Tiers}

To make deployment practical, we recommend three implementation tiers for Track B:

\begin{description}[leftmargin=1.5em]
  \item[Tier 1: Commitment and receipts.] Implement BN254 Pedersen commitments for Track B states and publish \(h_t^{(B)}\) as audit receipts. No encryption is required at this tier.
  \item[Tier 2: Transcript and multiplicity binding.] Add per-role transcript chains, context hash, and multiplicity operator \(M_t^{(B)}\). Use them to derive AEAD keys and to label Track B events.
  \item[Tier 3: Full QSBS and feedback dynamics.] Integrate joint frequency mapping, coupling tensor \(T_t^{(B)}\), and the feedback map \(f_B\) controlling adaptive key rotation and convergence-based policy decisions.
\end{description}

\subsection{Fastest Validation Path}

A practical validation plan for Track B is:
\begin{enumerate}[leftmargin=1.5em]
  \item Implement the BN254 Pedersen commitment and a minimal Track B wrapper that computes \(h_t^{(B)}\) from serialized state.
  \item Design a small Track B protocol (e.g., three-stage state update) and record transcripts, multiplicity vectors, and commitments.
  \item Verify that commitments recompute correctly and that multiplicity indices match manually counted message-class occurrences.
  \item Add transcript chains and context hash; confirm that modifying any message changes both \(\mathrm{ctx}_t^{(B)}\) and the derived Track B keys.
  \item Introduce a simple feedback rule for key rotation (e.g., increase rotation frequency when latency or error rises) and empirically confirm the convergence properties in a simulation.
\end{enumerate}

%-----------------------------------------------------------------------------%
\section{Defensive Publication Statement}
%-----------------------------------------------------------------------------%

\subsection*{Scope and Intent}

This document constitutes a formal defensive publication covering Track B as defined above, including:

\begin{itemize}[leftmargin=1.5em]
  \item The use of BN254 Pedersen commitments over multiplicity-encoded frequency representations of Track B state.
  \item The embedding of prime-indexed multiplicity operator states into HKDF info strings for directional AEAD key derivation.
  \item The design of Track B as a protocol whose transcript history is G\"odel-encoded via prime exponents, enabling recursive stability and audit.
  \item The integration of a rank-1 coupling tensor and feedback-driven contraction map for adaptive Track B key rotation and convergence.
\end{itemize}

\subsection*{Public-Domain Defensive Release}

By publishing this description with sufficient detail for an expert to re-implement and analyze the described mechanisms, the authors intend to establish prior art and prevent third parties from obtaining patent rights that would block or encumber implementations of Track B conforming to this specification.
This includes, but is not limited to, patent claims over the specific combination of multiplicity operators, BN254 Pedersen commitments, transcript-bound HKDF derivation, and adaptive feedback dynamics as applied to certification and secure transport in Track B.

\appendix
\section{Groups, Commitments, and Basic Bounds}

\subsection{BN254 and Pedersen Commitments}

Let \(\G_1\) be the prime-order subgroup of an elliptic curve (e.g., BN254) with prime order \(r\) and generator \(G\).
Let \(\F_r\) denote the associated scalar field.
We assume the standard hardness of the discrete logarithm problem in \(\G_1\).

\begin{definition}[Independent generator \(H\)]
Let \(\mathsf{HashToScalar} : \{0,1\}^* \to \F_r\) be a random-oracle hash.
Define
\[
  h_G = \mathsf{HashToScalar}(\texttt{"pedersen\_H"} \,\|\, \mathrm{enc}(G)) \in \F_r,
\]
and set
\[
  H = h_G \cdot G \in \G_1.
\]
\end{definition}

In the random-oracle model and under the discrete logarithm assumption, this gives a generator \(H\) that is computationally independent of \(G\).\cite{PedersenZKDocs,PedersenRareSkills,CommitmentSchemeWiki}

\begin{definition}[BN254 Pedersen commitment]
For value \(v \in \F_r\) and blinding factor \(\rho \in \F_r\), the Pedersen commitment is
\[
  C(v,\rho) = \rho G + v H \in \G_1.
\]
\end{definition}

\begin{proposition}[Additive homomorphism]
For any \(v_1,v_2,\rho_1,\rho_2 \in \F_r\),
\[
  C(v_1,\rho_1) + C(v_2,\rho_2) = C(v_1+v_2,\rho_1+\rho_2).
\]
\end{proposition}

\begin{proof}
Direct expansion:
\[
  C(v_1,\rho_1) + C(v_2,\rho_2) = (\rho_1 G + v_1 H) + (\rho_2 G + v_2 H)
  = (\rho_1+\rho_2)G + (v_1+v_2)H = C(v_1+v_2,\rho_1+\rho_2).
\]
This is the standard additive homomorphism of Pedersen commitments.\cite{PedersenZKDocs,PedersenRareSkills}
\end{proof}

\begin{proposition}[Hiding and binding (informal)]
Assuming the discrete logarithm problem in \(\G_1\) is hard and \(\rho\) is sampled uniformly at random from \(\F_r\), the scheme \(C(v,\rho)\) is:
\begin{itemize}
  \item Computationally \emph{hiding}: given \(C(v,\rho)\), distinguishing \(v\) from any other \(v'\) is infeasible.
  \item Computationally \emph{binding}: it is infeasible to find \((v,\rho) \ne (v',\rho')\) such that \(C(v,\rho) = C(v',\rho')\).
\end{itemize}
\end{proposition}

\begin{proof}[Proof sketch]
These properties are standard for Pedersen commitments in any group of prime order with independent generators.\cite{PedersenZKDocs,PedersenRareSkills,PedersenLibby,PedersenLibbitcoin}
Binding reduces to the hardness of discrete logarithm in \(\G_1\); hiding follows from the fact that \(\rho G\) is uniformly distributed in \(\G_1\) when \(\rho\) is uniform in \(\F_r\).
\end{proof}

\section{Multiplicity Operator and Transcript Encoding}

\subsection{Prime-Indexed Encoding}

Let \(\{p_i\}_{i \ge 1}\) be the increasing sequence of prime numbers.
Let \(\mathcal{C} = \{c_1,\dots,c_k\}\) denote the finite set of Track B message classes (e.g., \texttt{state\_update}, \texttt{ack}, \texttt{key\_rotate}, etc.).

\begin{definition}[Class counts]
Given a Track B transcript \(\sigma_t^{(B)}\) up to time \(t\), define
\[
  c_i(\sigma_t^{(B)}) \in \N
\]
as the number of messages of class \(c_i\) that appear in \(\sigma_t^{(B)}\).
\end{definition}

\begin{definition}[Prime-indexed transcript encoding]
The transcript encoding at time \(t\) is
\[
  \mathrm{enc}_B(\sigma_t^{(B)}) = \prod_{i=1}^k p_i^{c_i(\sigma_t^{(B)})} \in \N_{>0}.
\]
\end{definition}

\begin{definition}[Track B multiplicity operator]
The Track B multiplicity operator is the vector
\[
  M_t^{(B)} = \bigl(c_1(\sigma_t^{(B)}), \dots, c_k(\sigma_t^{(B)})\bigr) \in \N^k.
\]
\end{definition}

\begin{proposition}[Injectivity]
If \(\sigma_t^{(B)}\) and \(\sigma_t'{}^{(B)}\) have different multisets of message classes, then
\(\mathrm{enc}_B(\sigma_t^{(B)}) \ne \mathrm{enc}_B(\sigma_t'{}^{(B)})\).
Equivalently, \(\mathrm{enc}_B\) is injective with respect to the message-class multiset.
\end{proposition}

\begin{proof}
By the Fundamental Theorem of Arithmetic, the prime factorization of a positive integer is unique.
If two encodings were equal, the exponent of each prime \(p_i\) in the factorization would have to match, so \(c_i(\sigma_t^{(B)}) = c_i(\sigma_t'{}^{(B)})\) for all \(i\).
Thus the class multiplicities are equal and the multisets coincide.
\end{proof}

\subsection{Norm Bounds for \texorpdfstring{$M_t^{(B)}$}{M\_t\^B}}

Let
\[
  N_t = \sum_{i=1}^k c_i(\sigma_t^{(B)}) = \text{total number of messages in the transcript up to time } t.
\]

\begin{proposition}[Norm bounds on \(M_t^{(B)}\)]
For each time \(t\),
\begin{align*}
  \norm{M_t^{(B)}}_1 &= N_t, \\
  \norm{M_t^{(B)}}_2 &\le N_t, \\
  \norm{M_t^{(B)}}_\infty &\le N_t.
\end{align*}
\end{proposition}

\begin{proof}
The \(\ell_1\) norm is
\[
  \norm{M_t^{(B)}}_1 = \sum_{i=1}^k \abs{c_i(\sigma_t^{(B)})} = \sum_{i=1}^k c_i(\sigma_t^{(B)}) = N_t.
\]
For the \(\ell_\infty\) norm,
\[
  \norm{M_t^{(B)}}_\infty = \max_i c_i(\sigma_t^{(B)}) \le \sum_i c_i(\sigma_t^{(B)}) = N_t.
\]
For the \(\ell_2\) norm,
\[
  \norm{M_t^{(B)}}_2^2 = \sum_{i=1}^k c_i(\sigma_t^{(B)})^2 \le \max_i c_i(\sigma_t^{(B)}) \sum_i c_i(\sigma_t^{(B)}) \le N_t^2,
\]
so \(\norm{M_t^{(B)}}_2 \le N_t\).
\end{proof}

When used as a diagonal linear operator (see below), these bounds transfer directly to operator norms.

\subsection{Multiplicity as a Linear Operator}

To use \(M_t^{(B)}\) in operator-norm bounds, we view it as a diagonal matrix.

\begin{definition}[Diagonal multiplicity operator]
Given \(M_t^{(B)}\), define the diagonal matrix
\[
  \mathbf{M}_t^{(B)} = \mathrm{diag}\bigl(M_t^{(B)}\bigr) \in \R^{k \times k}.
\]
\end{definition}

\begin{proposition}[Operator norms of \(\mathbf{M}_t^{(B)}\)]
The spectral norm (operator 2-norm) of \(\mathbf{M}_t^{(B)}\) is
\[
  \norm{\mathbf{M}_t^{(B)}}_{2\to2} = \norm{M_t^{(B)}}_\infty \le N_t,
\]
and more generally \(\norm{\mathbf{M}_t^{(B)}}_{p\to p} = \norm{M_t^{(B)}}_\infty\) for any \(p \in [1,\infty]\).
\end{proposition}

\begin{proof}
For a diagonal matrix \(D = \mathrm{diag}(d_1,\dots,d_k)\), the operator norm \(\norm{D}_{p\to p}\) is \(\max_i \abs{d_i}\) for any \(p\).
This is immediate by considering unit vectors aligned with coordinates.
Thus
\[
  \norm{\mathbf{M}_t^{(B)}}_{2\to2} = \max_i \abs{c_i(\sigma_t^{(B)})} = \norm{M_t^{(B)}}_\infty \le N_t.
\]
\end{proof}

\section{Coupling Tensor and Encryption Operator}

\subsection{Rank-1 Coupling Tensor}

In the QSBS-enabled setting, the Track B state includes a joint representation of classical and quantum-like components, coupled via a rank-1 tensor.

Let \(u_t \in \R^n\) be the classical feature vector (e.g., frequency-coded) and \(v_t \in \R^m\) be the quantum-like amplitude feature vector for Track B at time \(t\).
Let their suprema over the time horizon be
\[
  B_c = \sup_t \norm{u_t}_2, \qquad
  B_q = \sup_t \norm{v_t}_2.
\]

\begin{definition}[Coupling tensor]
The Track B coupling tensor at time \(t\) is
\[
  T_t^{(B)} = u_t v_t^\top \in \R^{n \times m}.
\]
\end{definition}

\begin{proposition}[Norms of a rank-1 coupling tensor]
For \(T_t^{(B)} = u_t v_t^\top\),
\[
  \norm{T_t^{(B)}}_F = \norm{u_t}_2 \norm{v_t}_2, \qquad
  \norm{T_t^{(B)}}_{2\to2} = \norm{u_t}_2 \norm{v_t}_2,
\]
so in particular
\[
  \norm{T_t^{(B)}}_F \le B_c B_q, \qquad
  \norm{T_t^{(B)}}_{2\to2} \le B_c B_q.
\]
\end{proposition}

\begin{proof}
For a rank-1 matrix \(T = u v^\top\), the singular values are \(\{\norm{u}_2 \norm{v}_2,0,\dots,0\}\).
Hence the Frobenius norm satisfies
\[
  \norm{T}_F^2 = \sum_{i,j} u_i^2 v_j^2 = \norm{u}_2^2 \norm{v}_2^2,
\]
so \(\norm{T}_F = \norm{u}_2 \norm{v}_2\).
The spectral norm (largest singular value) is \(\norm{u}_2 \norm{v}_2\).
The inequalities follow by the definitions of \(B_c\) and \(B_q\).
\end{proof}

\subsection{Track B Encryption Operator}

Let \(S_t^{(B)}\) denote an embedded Track B state vector in \(\R^d\) for some fixed \(d\), such that \(T_t^{(B)} S_t^{(B)}\) and \(\mathbf{M}_t^{(B)} T_t^{(B)} S_t^{(B)}\) are well-defined.
Let \(F_B : \R^d \to \R^d\) be a (possibly nonlinear) feedback function.

\begin{definition}[Track B encryption operator]
The Track B encryption operator at time \(t\) is
\[
  E_B(t) = \mathbf{M}_t^{(B)} \bigl( T_t^{(B)} S_t^{(B)} \bigr) + F_B\bigl(S_t^{(B)}\bigr).
\]
\end{definition}

Assume there exists \(D_S, D_F > 0\) such that
\[
  \norm{S_t^{(B)}}_2 \le D_S, \qquad
  \norm{F_B(S_t^{(B)})}_2 \le D_F
\]
for all \(t\) in the horizon of interest.

\begin{proposition}[Norm bound on \(E_B(t)\)]
Under the assumptions above,
\[
  \norm{E_B(t)}_2 \le N_t B_c B_q D_S + D_F.
\]
\end{proposition}

\begin{proof}
Using the triangle inequality and submultiplicativity of operator norms,
\begin{align*}
  \norm{E_B(t)}_2
  &= \norm{\mathbf{M}_t^{(B)} T_t^{(B)} S_t^{(B)} + F_B(S_t^{(B)})}_2 \\
  &\le \norm{\mathbf{M}_t^{(B)} T_t^{(B)} S_t^{(B)}}_2 + \norm{F_B(S_t^{(B)})}_2 \\
  &\le \norm{\mathbf{M}_t^{(B)}}_{2\to2} \norm{T_t^{(B)}}_{2\to2} \norm{S_t^{(B)}}_2 + D_F.
\end{align*}
By the earlier bounds,
\[
  \norm{\mathbf{M}_t^{(B)}}_{2\to2} \le N_t, \quad
  \norm{T_t^{(B)}}_{2\to2} \le B_c B_q, \quad
  \norm{S_t^{(B)}}_2 \le D_S,
\]
so
\[
  \norm{E_B(t)}_2 \le N_t B_c B_q D_S + D_F.
\]
\end{proof}

\section{HKDF, Transcript Binding, and Multiplicity Labels}

\subsection{Transcript Chains and Context Hash}

Let \(\SHA\) be a hash function modeled as a random oracle.
For each role \(s \in \{A,B\}\), define a per-role transcript chain:

\begin{definition}[Per-role transcript chain]
Initialize
\[
  \mathrm{chain}_{s,0} = 0^{256}.
\]
For \(t \ge 1\), given the \(t\)-th message \(m_{s,t}\) sent by role \(s\),
\[
  \mathrm{chain}_{s,t} = \SHA\bigl( \mathrm{chain}_{s,t-1} \,\|\, m_{s,t} \bigr).
\]
\end{definition}

\begin{definition}[Track B context hash]
The Track B context hash at time \(t\) is
\[
  \mathrm{ctx}_t^{(B)} = \SHA\bigl( \mathrm{chain}_{A,t} \,\|\, \mathrm{chain}_{B,t} \bigr).
\]
\end{definition}

These definitions mirror standard hash-chain transcript binding patterns used in secure messaging and zero-knowledge transcript binding.\cite{CommitmentSchemeWiki}

\subsection{Multiplicity Label Encoding}

\begin{definition}[Multiplicity label encoding]
Given multiplicity vector \(M_t^{(B)}\), define a fixed-length encoding
\[
  \mathrm{mult}_t^{(B)} = \mathrm{Enc}_8\bigl(M_t^{(B)}\bigr) \in \{0,1\}^{64},
\]
where \(\mathrm{Enc}_8\) is an injective (or collision-resistant) map from the relevant multiplicity range into 8 bytes.
\end{definition}

This label is treated as authenticated auxiliary information in the key-derivation process, not as a source of cryptographic entropy.

\subsection{HKDF-Based Directional Key Derivation}

Let \(\HKDF\) denote an HMAC-based key-derivation function following the extract-then-expand paradigm (e.g., RFC 5869).\cite{HKDFWiki,HKDFDotNet,GopherHKDF}

\begin{definition}[Directional info string]
For direction \(d \in \{A \to B, B \to A\}\),
\[
  \mathrm{info}_d^{(B)} =
    \texttt{"QSBS-v1:enc:"} \,\|\, d \,\|\, \mathrm{ctx}_t^{(B)} \,\|\, \mathrm{mult}_t^{(B)}.
\]
\end{definition}

\begin{definition}[Directional AEAD key for Track B]
Let \(K_t\) be a shared secret (e.g., from Diffie–Hellman or QKD) and \(\mathrm{salt}_t\) a salt.
The directional AEAD key at time \(t\) is
\[
  K_{\mathrm{enc},d,t}^{(B)} = \HKDF\bigl(K_t; \mathrm{salt}_t, \mathrm{info}_d^{(B)}\bigr).
\]
\end{definition}

\begin{remark}
Standard analyses of HKDF show that if \(K_t\) has sufficient min-entropy and HMAC behaves as a pseudorandom function, then HKDF outputs are indistinguishable from random keys.\cite{HKDFWiki,HKDFDotNet,GopherHKDF}
Including \(\mathrm{ctx}_t^{(B)}\) and \(\mathrm{mult}_t^{(B)}\) in \(\mathrm{info}_d^{(B)}\) binds the derived keys to a specific transcript and multiplicity label without degrading security, provided these values are public and not reused inconsistently.
\end{remark}

\section{AEAD Encryption and Associated Data}

Let \(\Enc_K(\mathrm{nonce}, m, \mathrm{ad})\) be an AEAD encryption algorithm, such as AES-GCM, with key \(K\), nonce, plaintext \(m\), and associated data \(\mathrm{ad}\).
Assume the AEAD is IND-CPA and INT-CTXT secure.

\begin{definition}[Track B associated data]
Let \(x_t^{(B)}\) be the Track B plaintext payload at time \(t\), and let \(F_t^{(B)}\) and \(M_t^{(B)}\) be the associated frequency vector and multiplicity operator.
Define associated data as the serialized tuple
\[
  \mathrm{ad}_t^{(B)} = \mathrm{enc}\bigl(\mathrm{ctx}_t^{(B)}, F_t^{(B)}, M_t^{(B)}\bigr).
\]
\end{definition}

\begin{definition}[Track B ciphertext]
Let \(\mathrm{nonce}_{d,t}\) be a nonce.
The Track B ciphertext for direction \(d\) at time \(t\) is
\[
  C_t^{(B)} = \Enc_{K_{\mathrm{enc},d,t}^{(B)}}\bigl(\mathrm{nonce}_{d,t}, x_t^{(B)}, \mathrm{ad}_t^{(B)}\bigr).
\]
\end{definition}

\begin{remark}
Because AEAD schemes authenticate the associated data, any attempt to tamper with \(\mathrm{ctx}_t^{(B)}\), \(F_t^{(B)}\), or \(M_t^{(B)}\) will lead to decryption failure.
This provides integrity and binding between the ciphertext, the transcript state, and the multiplicity operator.
\end{remark}

\section{Feedback Dynamics and Contraction}

\subsection{State Space and Metric}

Let the Track B feedback state be a pair
\[
  Z_t^{(B)} = \bigl(M_t^{(B)}, T_t^{(B)}\bigr) \in \R^k \times \R^{n \times m}.
\]
Equip this product space with the norm
\[
  \norm{Z_t^{(B)}} := \norm{M_t^{(B)}}_2 + \norm{T_t^{(B)}}_F.
\]

\begin{lemma}[Completeness]
The space \(\R^k \times \R^{n \times m}\) with the metric \(d(Z,Z') = \norm{Z-Z'}\) is a complete metric space.
\end{lemma}

\begin{proof}
Both \(\R^k\) and \(\R^{n \times m}\) are finite-dimensional normed spaces and hence complete; their Cartesian product equipped with the sum norm is also complete.
\end{proof}

\subsection{Feedback Map and Lipschitz Condition}

Let \(\Theta_t\) denote a vector of observed metrics at time \(t\) (e.g., error rates, latency, QBER, load).
Define a feedback map \(f_B\) via
\[
  Z_{t+1}^{(B)} = f_B\bigl(Z_t^{(B)}, \Theta_t\bigr).
\]

\begin{definition}[Lipschitz condition]
Suppose there exist non-negative constants \(L_M,L_T\) such that for all \(Z = (M,T)\) and \(Z' = (M',T')\),
\begin{align*}
  \norm{f_B(Z,\Theta)_M - f_B(Z',\Theta)_M}_2 &\le L_M \norm{Z-Z'}, \\
  \norm{f_B(Z,\Theta)_T - f_B(Z',\Theta)_T}_F &\le L_T \norm{Z-Z'}.
\end{align*}
Let \(L = L_M + L_T\).
\end{definition}

\begin{proposition}[Global Lipschitz bound]
Under the above conditions,
\[
  \norm{f_B(Z,\Theta) - f_B(Z',\Theta)} \le L \norm{Z - Z'}
\]
for all \(Z,Z'\).
\end{proposition}

\begin{proof}
By definition,
\[
  \norm{f_B(Z,\Theta) - f_B(Z',\Theta)}
  = \norm{f_B(Z,\Theta)_M - f_B(Z',\Theta)_M}_2
    + \norm{f_B(Z,\Theta)_T - f_B(Z',\Theta)_T}_F.
\]
Applying the component-wise Lipschitz bounds yields
\[
  \norm{f_B(Z,\Theta) - f_B(Z',\Theta)}
  \le (L_M + L_T) \norm{Z-Z'} = L \norm{Z-Z'}.
\]
\end{proof}

\begin{theorem}[Banach contraction for Track B feedback]
If \(L < 1\), then:
\begin{enumerate}
  \item There exists a unique fixed point \(Z_\ast^{(B)}\) such that
  \[
    f_B\bigl(Z_\ast^{(B)},\Theta\bigr) = Z_\ast^{(B)}
  \]
  for all admissible \(\Theta\) in the regime considered (e.g., steady-state metrics).
  \item For any initial state \(Z_0^{(B)}\), the iterations
  \[
    Z_{t+1}^{(B)} = f_B\bigl(Z_t^{(B)},\Theta_t\bigr)
  \]
  satisfy the contraction bound
  \[
    \norm{Z_t^{(B)} - Z_\ast^{(B)}} \le L^t \norm{Z_0^{(B)} - Z_\ast^{(B)}},
  \]
  provided the Lipschitz constants remain valid over the trajectory.
\end{enumerate}
\end{theorem}

\begin{proof}
By the previous lemma, the state space is complete.
If \(L<1\), then for fixed \(\Theta\) (or for \(\Theta_t\) staying in a regime where the same Lipschitz constants apply), the map \(Z \mapsto f_B(Z,\Theta)\) is a contraction.
The Banach Fixed-Point Theorem guarantees existence and uniqueness of \(Z_\ast^{(B)}\) and geometric convergence from any starting point.
\end{proof}

\section{Summary of Bounds and Guarantees}

For Track B, the key operator norm and security-relevant bounds can be summarized as follows:

\begin{itemize}
  \item \textbf{Multiplicity norms.}
  For \(M_t^{(B)} \in \N^k\),
  \[
    \norm{M_t^{(B)}}_1 = N_t, \quad
    \norm{M_t^{(B)}}_2 \le N_t, \quad
    \norm{M_t^{(B)}}_\infty \le N_t,
  \]
  and for the diagonal operator \(\mathbf{M}_t^{(B)}\),
  \(\norm{\mathbf{M}_t^{(B)}}_{2\to2} \le N_t\).
  \item \textbf{Coupling tensor.}
  For the rank-1 tensor \(T_t^{(B)} = u_t v_t^\top\),
  \[
    \norm{T_t^{(B)}}_F = \norm{u_t}_2 \norm{v_t}_2 \le B_c B_q, \quad
    \norm{T_t^{(B)}}_{2\to2} \le B_c B_q.
  \]
  \item \textbf{Encryption operator.}
  For bounded state \(\norm{S_t^{(B)}}_2 \le D_S\) and feedback \(\norm{F_B(S_t^{(B)})}_2 \le D_F\),
  \[
    \norm{E_B(t)}_2 \le N_t B_c B_q D_S + D_F.
  \]
  \item \textbf{HKDF security.}
  Provided \(K_t\) has sufficient min-entropy and HKDF is instantiated with a secure HMAC, \(K_{\mathrm{enc},d,t}^{(B)}\) is computationally indistinguishable from random, even when \(\mathrm{info}_d^{(B)}\) includes public multiplicity labels and context hashes.\cite{HKDFWiki,GopherHKDF}
  \item \textbf{AEAD binding.}
  AEAD security ensures integrity of both ciphertext and associated data; any change in \(\mathrm{ctx}_t^{(B)}\), \(F_t^{(B)}\), or \(M_t^{(B)}\) will cause decryption to fail, binding Track B ciphertexts to their specific transcript and multiplicity state.
  \item \textbf{Feedback contraction.}
  If the Track B feedback map satisfies the Lipschitz condition with \(L<1\), the multiplicity and coupling tensor converge to a unique fixed point at geometric rate, providing a mathematically controlled regime for adaptive key rotation or state evolution.
\end{itemize}

These results form the mathematical backbone of the Track B design and justify the use of prime-indexed multiplicities, rank-1 coupling tensors, and HKDF-based directional keys in the combined certification-and-transport architecture.

\vspace{1em}
\noindent\textbf{References (for background only)}
\begin{itemize}
  \item Standard references on Walsh--Hadamard transforms and orthogonal transforms.\cite{WHTOverview,WHTOrthogonal}
  \item Cryptographic references on Pedersen commitments, HKDF, and commitment schemes.\cite{PedersenZKDocs,PedersenRareSkills,CommitmentSchemeWiki,HKDFWiki,GopherHKDF}
\end{itemize}

\nocite{*}
\bibliographystyle{unsrt}  
\bibliography{references}  


\end{document}
